% EGA V, English translation by P. Blass, 1987 (cote 24): archive cote 24, site page 74 (PDF page 75).
% typescript. Draft compared: none.
% AI-transcribed by Claude Opus 5.5 (claude-opus-5-5, adaptive thinking, effort high), Message Batches API, batch ega24b (msgbatch_01RGJJeHdR96U16ZWFuUjtDx), 2026-11-13; not yet checked by eye.
\scanpage{24/74}
\mspage{74}

say) as linear subvarieties of codim 1 of P.  (called \uline{again}
\uline{hyperplanes} of P).

This says that $L^0$ consists of hyperplanes which \uline{contain} the
linear subvariety L of P [by which of course we mean] that the
points of L with values in $S^1$ are the hyperplanes in $P_S{}^1$ that
contain $L_S{}^1$).  This follows from the fact that should have occured
at the same time as the fact that a linear subvariety L of $\underline{P}$
determines a locally free quotient G or Q [illegible] of E from
which it results that if G and $G^1$ are two locally free quotients of
E (not necessarily of the same rank) then $\underline{P}(G^1)$ C $\underline{P}(G)$ in so far as
the linear subvarieties of $\underline{P}(E)$, $G^1$ is majorized by G (and the
inclusion $\underline{P}(G^1) \rightarrow \underline{P}(G)$ is nothing else but the deduced morphism G
$\rightarrow G^1$).

Here is a minimum of the [sorites Fr]* which we must have at
our disposal.  The complete list cannot in any case be fixed (only
when) until the sets of other N$^{o1}$s of the present paragraph are
written up.

It seems to me convenient [commode Fr] to introduce also the
functor Grass (E) ($S^1$) = set of quotient modules locally free (of
rank not specified) of $E_S1$ [illegible ask AG]
then \uline{Grass} (E) is representable by n $\geq$ 0 $\mathrm{Grass}_n$(E).  The linear
subvarieties of $\underline{P}$(E) are indeed defined by sections of Grass (E)
over S [NB:  the rank i.e. the relative dimension may vary if S is
not connected.  [slightly illegible confirm with AG].

\margin{* Ask AG or Deligne about the best word.}
